\documentclass{article}
\usepackage[spanish]{babel}
\usepackage[utf8]{inputenc}
\usepackage{amsmath}
\usepackage{booktabs}
\usepackage{array}
\usepackage{geometry}
\usepackage{circuitikz}
\geometry{a4paper, margin=1in}

\title{Prácticas de Electricidad - Teorema de Thévenin}
\author{}
\date{}

\begin{document}

\maketitle

\section*{Página 156 - Práctica de Teorema de Thévenin}

\begin{table}[h]
\centering
\caption{TABLA 26-1}
\label{tab:26-1}
\begin{tabular}{|c|c|c|c|c|c|c|c|}
\hline
$V_{TH}$ (voltios) & $R_{TH}$ (ohmios) & \multicolumn{2}{c|}{$I_L$ (mA)} \\
\cline{3-4}
& & \multicolumn{2}{c|}{} \\
\hline
$R_L$ & Medida & Calculada & Medida & Calculada & Medida (circuito original) & Medida (Thévenin) & Calculada \\
\hline
2.700 & & & & & & & \\
\hline
& & & & & & & \\
\hline
& & & & & & & \\
\hline
\end{tabular}
\end{table}

tro M entre los puntos C y B, como en la figura 26-4.

\begin{enumerate}
\setcounter{enumi}{4}
\item Cerrar el interruptor S. Vigilar $V_{TH}$ para que no varíe su valor. Medir la corriente en $R_L$ y anotarla en la columna «$I_L$ medida Thévenin».

\item Desconectar la alimentación. En el circuito de la figura 26-3 calcular el valor de $V_{TH}$ para el generador equivalente de Thévenin y la carga $R_L$ y anotarla en la tabla 26-1 en la columna «$V_{TH}$ calculada». Presentar los cálculos. Calcular también el valor de $R_{TH}$ y anotarlo en la columna «$R_{TH}$ calculada». Presentar los cálculos. Hallar ahora los valores de $I_L$ en $R_L$ haciendo uso del circuito equivalente de Thévenin con los valores calculados. Presentar los cálculos. Anotar en tabla 26-1 en columna «$I_L$ calculada».

\textbf{Adicional}
\item Comprobar que es cierto el teorema de Thévenin para el circuito de la figura 26-3 con otros valores de $R_L$, sustituyendo sucesivamente ésta por tres resistencias cuyos valores no sean menores de 1.200 ohmios ni mayores de 10.000 ohmios. Medir la corriente $I_L$ en $R_L$ en el circuito original y también en el circuito equivalente de Thévenin. Anotar en tabla 26-1 los valores de $R_L$ empleados y los otros datos pertinentes. Explicar detalladamente el método seguido.
\end{enumerate}

\begin{figure}[h!]
    \centering
    \begin{circuitikz}[scale=0.9, transform shape]
        \shorthandoff{<}\shorthandoff{>}
        % Definir coordenadas
        \coordinate (P) at (0, 4);
        \coordinate (B) at (5, 4);
        \coordinate (C) at (5, 0);
        \coordinate (S_node) at (2, 0);
        \coordinate (V_neg) at (0, 0);

        % Dibujar componentes
        \draw (P) to[R, l={$5,6\,\text{k}\Omega$}] ++(2.5,0) to[R, l={$10\,\text{k}\Omega$}] (B);
        \draw (B) to[R, l_=$R_L$, l^=$2.700$] (B |- C) node[midway, right, xshift=2mm] {} to[ammeter, l=M, a^=+, a_=-] (C);
        \draw (C) to[short] (S_node);
        \draw (S_node) to[closing switch, l=S] (V_neg);
        \draw (V_neg) to[V, l_=$V_{TH}$, v_<={}] (P);

        % Etiquetas de nodos
        \node[above left] at (P) {P};
        \node[above right] at (B) {B};
        \node[below right] at (C) {C};
    \end{circuitikz}
    \caption{FIG. 26-4. — Equivalente experimental de Thévenin correspondiente a la figura 26-3.}
    \label{fig:26-4}
\end{figure}

\section*{PREGUNTAS}

\begin{enumerate}
\item ¿Cómo es el valor de $I_L$ medido en el circuito de la figura 26-3 comparado con el del circuito equivalente de Thévenin de la figura 26-4? ¿Deben ser iguales? Explicarlo.

\item Discutir los datos de la tabla 26-1 y explicar los resultados inesperados que pudieran encontrarse.

\item ¿Confirma el experimento la validez del teorema de Thévenin para el circuito de la figura 26-3? Explicarlo.

\item ¿«Prueba» este experimento el teorema de Thévenin para cualquier circuito? Explicarlo.

\item ¿Qué ventajas específicas, si las hay, cree Ud. que tiene el uso del teorema de Thévenin en la solución de circuitos de c.c.?
\end{enumerate}

\section*{Respuestas a las preguntas de auto-examen}

(Con aproximación de la regla de cálculo).

\begin{enumerate}
\item \begin{enumerate}
    \item 100 V
    \item 188 $\Omega$
    \item 0,134 A
    \end{enumerate}
\item \begin{enumerate}
    \item 45 V
    \item 338 $\Omega$
    \item 0,1 A
    \end{enumerate}
\end{enumerate}

\end{document}